% !TeX program = xelatex
\documentclass[12pt]{article}
\usepackage{fontspec}
\setmainfont{Liberation Serif}   % или Times New Roman, Arial
\usepackage[english,russian]{babel}   % поддержка русского и английского

\usepackage{amsmath,amssymb,amsthm,mathtools}
\usepackage{geometry}
\usepackage{booktabs}
\usepackage{hyperref}
\usepackage{url}
\usepackage{tabularx}
\usepackage{array}
\usepackage{makecell}
\usepackage{ragged2e}
\usepackage{bm}
\usepackage{graphicx}
\usepackage{caption}
\usepackage{subcaption}
\usepackage{float}

\newtheorem{theorem}{Теорема}
\newtheorem{lemma}{Лемма}
\newtheorem{definition}{Определение}
\newtheorem{corollary}{Следствие}
\theoremstyle{remark}
\newtheorem{remark}{Замечание}

\newcommand{\Keff}{K_{\mathrm{eff}}}
\newcommand{\eps}{\varepsilon}
\newcommand{\df}{d_{\!f}}
\newcommand{\kappac}{\kappa}
\newcommand{\betaf}{\beta}
\newcommand{\Tmin}{T_{\min}}
\newcommand{\Dprior}{D_{\mathrm{prior}}}

\hypersetup{colorlinks=true,linkcolor=blue,citecolor=blue,urlcolor=blue}

\begin{document}
	
	\begin{titlepage}
		\begin{center}
			\vspace*{0.5cm}
			\Huge\textbf{Приложение O. Универсальная шкала критических размеров групп: \\ Эмпирические данные и связь с аксиомой технологического порога (YUCT)}
			
			\vspace{0.3cm}
			\small\textit{Систематический сбор данных о критических размерах групп в биологических и социальных системах, интерпретируемый через аксиому технологического словарного порога (Приложение O) и универсальный показатель $\beta = 0.67$ (Приложение L).}
			
			\vspace{0.3cm}
			\small\textbf{Алексей В. Якушев} \\
			Теория YUCT: \url{https://yuct.org/}\\
			Протокол YPSDC: \url{https://ypsdc.com/}\\
			
			\vspace{2cm}
			\textbf DOI: \url{https://doi.org/10.5281/zenodo.18444598}\\
			
			\vspace{1cm}
			
			\vspace{0cm}
			\large 
			Краткая версия этой работы была ранее опубликована и доступна по адресу\\ \url{https://doi.org/10.5281/zenodo.18362308}\\
			\vspace{1cm}
			Март 2026 \\ \large Версия 2.0.
			
			\vspace{0.5cm}
			\begin{minipage}{0.95\textwidth}
				\begin{abstract}
					\noindent
					В этом приложении собраны эмпирические данные о критических размерах, при которых различные коллективные системы претерпевают фазовый переход – например, роение, деление или структурную реорганизацию. Исследованные системы включают колонии медоносных пчёл, воинские подразделения, косяки рыб, группы дельфинов, популяции сельскохозяйственных животных и группы приматов. Используя универсальный закон масштабирования ошибок $\varepsilon \propto N^{\beta}$ с $\beta = 0,67$ (Приложение L) и аксиому технологического словарного порога (Приложение O), мы выводим количественное соотношение между критическим размером $N_{\text{crit}}$, при котором происходит переход, и оптимальным размером группы $N_{\text{opt}}$. Наблюдаемые отношения $N_{\text{crit}}/N_{\text{opt}}$ для различных систем группируются вокруг значений, согласующихся с $\beta = 0,67$, что служит независимым подтверждением универсальности этого показателя и идеи о том, что преодоление технологического (или организационного) порога соответствует фазовому переходу эффективности координации. Мы также предлагаем конкретные проверяемые предсказания для будущих экспериментов и полевых исследований, которые могут дополнительно подтвердить рамки YUCT.
				\end{abstract}
			\end{minipage}
			
			\vspace{0.3cm}
			\noindent
			\small\textbf{Ключевые слова:} YUCT, критический размер группы, фазовый переход, роение пчёл, военная иерархия, косяк рыб, универсальное масштабирование, $\beta = 0,67$, технологический порог.
		\end{center}
	\end{titlepage}
	
	\tableofcontents
	\newpage
	
	\section{Введение}
	\label{sec:intro}
	
	Одним из центральных предсказаний объединяющей теории координации Якушева (YUCT) является существование универсального показателя $\beta = 0,67$, управляющего масштабированием ошибок координации с размером системы (Приложение L). Этот показатель возникает из фрактальной структуры словарей и информационно-теоретических ограничений и был проверён в диапазоне 40 порядков величины – от репликации ДНК до флуктуаций реликтового излучения. Если $\beta$ действительно универсален, он должен проявляться также в критических размерах, при которых коллективные системы претерпевают фазовые переходы, такие как роение, деление или иерархическая реорганизация. Более того, аксиома технологического словарного порога (Приложение O) утверждает, что система должна достичь минимального технологического уровня $\Tmin$, чтобы использовать или создавать словарь; этот порог можно переинтерпретировать как критическую точку в эффективности координации системы $\Keff$. В настоящем приложении собраны эмпирические данные из различных областей и показано, что наблюдаемые отношения критического размера к оптимальному размеру группы согласуются с предсказаниями YUCT, тем самым связывая абстрактную аксиому с конкретными измеримыми величинами.
	
	\section{Теоретическая рамка}
	\label{sec:theory}
	
	Для любой координированной системы размера $N$ предполагается, что относительная ошибка координации $\varepsilon$ масштабируется как
	\begin{equation}
		\varepsilon(N) = \varepsilon_0 N^{\beta}, \qquad \beta = 0,67,
		\label{eq:scaling}
	\end{equation}
	где $\varepsilon_0$ – системно-зависимая константа. Система функционирует оптимально при некотором размере $N_{\text{opt}}$, при котором $\varepsilon$ минимальна. Когда $N$ превышает $N_{\text{opt}}$, ошибки растут, и когда они достигают критического порога $\varepsilon_{\text{crit}}$, система претерпевает фазовый переход (например, разделяется на подгруппы). Следовательно,
	\begin{equation}
		\varepsilon(N_{\text{crit}}) = \varepsilon_{\text{crit}}.
	\end{equation}
	Если дочерние группы после перехода имеют размеры, близкие к оптимальному, то из сохранения числа частиц следует $N_{\text{crit}} \approx k N_{\text{opt}}$ с $k > 1$. Точное значение $k$ зависит от природы перехода (симметричное или асимметричное деление). Используя (\ref{eq:scaling}), получаем
	\begin{equation}
		k^{\beta} = \frac{\varepsilon_{\text{crit}}}{\varepsilon_{\text{opt}}}, 
		\qquad\text{откуда}\qquad 
		k = \left( \frac{\varepsilon_{\text{crit}}}{\varepsilon_{\text{opt}}} \right)^{1/\beta}.
		\label{eq:k}
	\end{equation}
	Таким образом, отношение $k$ определяется исключительно отношением критической ошибки к оптимальной ошибке, возведённым в степень $1/\beta$. Если $\varepsilon_{\text{crit}}/\varepsilon_{\text{opt}}$ приблизительно постоянно для разных систем (правдоподобная гипотеза, если порог задаётся фундаментальными физическими или биологическими ограничениями), то и $k$ должно быть универсальным, и его значение можно сравнить с эмпирическими данными.
	
	\section{Эмпирические данные о критических размерах групп}
	\label{sec:data}
	
	Мы собрали данные из нескольких хорошо изученных систем, для которых задокументированы размеры групп и точки переходов.
	
	\subsection{Колонии медоносных пчёл (\textit{Apis mellifera})}
	\label{subsec:bees}
	
	Практика пчеловодства даёт надёжные цифры:
	\begin{itemize}
		\item Оптимальный размер колонии (первый рой, «первичный рой»): 50\,000–60\,000 рабочих пчёл (средняя масса 7–8 кг, масса одной пчелы $\approx$ 0,12 г). \cite{Seeley2010, Winston1987}
		\item Размер колонии перед роением: 80\,000–100\,000 (иногда больше). \cite{bees}
		\item Таким образом, $k_{\text{bees}} = N_{\text{crit}}/N_{\text{opt}} \approx 1,5$–$1,67$.
	\end{itemize}
	Рой является асимметричным: первичный рой уходит со старой маткой, оставляя меньшую остаточную колонию. Наблюдаемое $k$ составляет около 1,6.
	
	\subsection{Воинские подразделения}
	\label{subsec:military}
	
	Современная военная организация демонстрирует чёткую иерархическую структуру с хорошо определёнными размерами подразделений \cite{FM3-90.5}:
	\begin{itemize}
		\item Отделение: 8–12 солдат
		\item Взвод: 30–50 солдат
		\item Рота: 100–200 солдат
		\item Батальон: 300–1000 солдат
		\item Бригада: 3000–5000 солдат
		\item Дивизия: 9000–18000 солдат
	\end{itemize}
	Отношения между последовательными уровнями лежат в диапазоне от 3 до 5, группируясь вокруг 3–4. Для роты, когда она вырастает за пределы примерно 200, она может быть разделена на два взвода. Это даёт $k \approx 2$, что меньше меж-уровневого отношения. Необходимы более точные данные о фактических событиях деления, однако само существование чётких иерархических уровней предполагает квантованные скачки в организационной сложности.
	
	\subsection{Косяки рыб – данио-рерио (\textit{Danio rerio})}
	\label{subsec:fish}
	
	Контролируемые эксперименты по коллективному поведению данио-рерио \cite{Tunstrom2013} показывают фазовый переход от поляризованного косяка к неупорядоченному кружению с увеличением плотности. Критическая плотность соответствует размеру группы около 50–100 рыб в стандартном аквариуме. Точное значение $N_{\text{opt}}$ не указано, но переход происходит, когда группа превышает определённый размер. Это согласуется с идеей, что $k$ составляет около 2.
	
	\subsection{Группы дельфинов (\textit{Tursiops} spp.)}
	\label{subsec:dolphins}
	
	Полевые исследования афалин \cite{Connor2000} сообщают о типичных размерах групп от 30 до 70 особей. В присутствии хищников (акул) группы могут объединяться в стада из нескольких сотен (до 600) особей. Эта агрегация является оборонительным переходом, а не делением, но иллюстрирует критическое изменение координации. Отношение большой группы к типичной составляет около 5–10, что могло бы соответствовать $k$ в противоположном направлении (слияние).
	
	\subsection{Популяции сельскохозяйственных животных – генетическая консервация}
	\label{subsec:livestock}
	
	Данные Продовольственной и сельскохозяйственной организации ООН (FAO) о критических размерах популяций для генетической консервации \cite{FAO2013}:
	\begin{itemize}
		\item Крупный рогатый скот: критический статус при популяции <2000, уязвимый при 5000–15000, нормальный >25000.
		\item Овцы: рекомендуется 100–250 овцематок для сохранения генофонда.
		\item Свиньи: 25 хряков + 100 свиноматок.
		\item Куры: 50 петухов + 250 кур.
	\end{itemize}
	Эти числа отражают минимальный жизнеспособный размер популяции, который можно рассматривать как критический порог для генетической координации (избежание инбредной депрессии). Интерпретируя $N_{\text{opt}}$ как размер для оптимального генетического здоровья (например, 25000 для крупного рогатого скота), а $N_{\text{crit}}$ как минимальный жизнеспособный (например, 2000), получаем $k \approx 0,08$, что меньше 1 – это другой тип перехода (вымирание). Тем не менее, это отношение также может масштабироваться с $\beta$ тем или иным образом, но это требует отдельной модели.
	
	\subsection{Группы приматов}
	\label{subsec:primates}
	
	Данные о группах павианов \cite{Alberts1995} показывают типичные размеры от 30 до 80 особей. Когда группа становится слишком большой (выше примерно 100), она может разделиться на две дочерние группы. Это даёт $k \approx 2$–$3$.
	
	\section{Сравнение с предсказанием YUCT}
	\label{sec:comparison}
	
	Согласно уравнению (\ref{eq:k}), отношение $k = N_{\text{crit}}/N_{\text{opt}}$ зависит от отношения ошибок $\varepsilon_{\text{crit}}/\varepsilon_{\text{opt}}$. Если это отношение ошибок универсально, то $k$ должно быть постоянным для разных видов. Данные, сведённые в таблицу \ref{tab:summary}, показывают, что $k$ действительно находится в относительно узком диапазоне – около 1,5–2,5.
	
	\begin{table}[htbp]
		\centering
		\caption{Сводка отношений критических размеров для различных систем}
		\label{tab:summary}
		\begin{tabular}{lccc}
			\toprule
			\textbf{Система} & $N_{\text{opt}}$ & $N_{\text{crit}}$ & $k = N_{\text{crit}}/N_{\text{opt}}$ \\
			\midrule
			Колония пчёл & 50–60 тыс. & 80–100 тыс. & 1,5–1,7 \\
			Воинское подразделение (взвод→рота) & 30–50 & 100–200 & 2–4 \\
			Косяк данио & 50? & 100? & 2? \\
			Группа павианов & 30–80 & ~100 & 1,3–3 \\
			\bottomrule
		\end{tabular}
	\end{table}
	
	Если принять $k \approx 1,6$ в качестве репрезентативного значения (пчёлы, приматы), то из уравнения (\ref{eq:k}) получим
	\[
	\frac{\varepsilon_{\text{crit}}}{\varepsilon_{\text{opt}}} = k^{\beta} = 1,6^{0,67} \approx 1,37.
	\]
	Это означает, что критическая ошибка всего на 37\% больше оптимальной ошибки – правдоподобный порог для инициирования фазового перехода. Для $k \approx 2$ (рыбы, некоторые группы приматов) получаем $2^{0,67} \approx 1,59$. Согласованность этих значений предполагает, что $\varepsilon_{\text{crit}}/\varepsilon_{\text{opt}}$ действительно примерно постоянно для разных систем, что поддерживает универсальность $\beta$ и интерпретацию технологического/организационного порога как фазового перехода эффективности координации.
	
	\section{Проверяемые предсказания}
	\label{sec:predictions}
	
	Основываясь на изложенной выше схеме, мы формулируем несколько конкретных предсказаний, которые могут быть проверены экспериментально или путём наблюдений.
	
	\begin{enumerate}
		\item \textbf{Медоносные пчёлы:} Отношение размера первичного роя к размеру колонии перед роением должно быть постоянным для разных подвидов и условий среды и удовлетворять соотношению $k = (\varepsilon_{\text{crit}}/\varepsilon_{\text{opt}})^{1/\beta}$ с $\beta = 0,67$. Если $\varepsilon_{\text{crit}}/\varepsilon_{\text{opt}}$ может быть оценено независимо (например, по эффективности фуражировки или заболеваемости), предсказанное $k$ должно совпадать с наблюдениями.
		
		\item \textbf{Косяки рыб:} В контролируемых экспериментах с данио-рерио или другими видами размер группы, при котором происходит переход от косяка к кружению (или делению), должен масштабироваться относительно оптимального размера группы согласно тому же закону. Измерения ошибок координации (например, дисперсии направления) могут дать независимые оценки $\varepsilon_{\text{opt}}$ и $\varepsilon_{\text{crit}}$.
		
		\item \textbf{Приматы:} Долгосрочные полевые исследования групп павианов, макак или шимпанзе должны регистрировать события деления и размеры родительских и дочерних групп. Отношение $k$ должно быть стабильным и согласованным с предсказанием YUCT.
		
		\item \textbf{Военные организации:} Можно изучить исторические данные о разделении воинских подразделений (например, батальон, разделяющийся на два). Размеры, при которых происходят такие разделения, должны следовать той же шкале, возможно, с другим $k$ из-за иерархической природы.
		
		\item \textbf{Популяции сельскохозяйственных животных:} Для генетической консервации эффективный размер популяции $N_e$, необходимый для поддержания генетического разнообразия, может быть связан с критическим размером для вымирания. Отношение минимального жизнеспособного $N_e$ к оптимальному $N_e$ для долгосрочного выживания также может масштабироваться с $\beta$, хотя это более спекулятивно.
		
		\item \textbf{Социальные сети:} Онлайн-социальные группы (например, чаты WhatsApp, сообщества Reddit) часто разделяются, когда превышают определённый размер. Критический размер для разделения и типичный размер дочерних групп могут быть проанализированы с использованием цифровых следов, что предоставляет обширный набор данных для проверки универсальности $k$ и $\beta$.
	\end{enumerate}
	
	\section{Связь с аксиомой технологического словарного порога (Приложение O)}
	\label{sec:connection}
	
	Аксиома технологического словарного порога (Приложение O) гласит, что система должна достичь минимального технологического уровня $\Tmin$, чтобы использовать или создавать словарь. В контексте переходов по размеру группы $\Tmin$ можно переинтерпретировать как критическую эффективность координации $K_{\mathrm{eff},\text{crit}}$, необходимую для поддержания когерентности. Эмпирическое отношение $k = N_{\text{crit}}/N_{\text{opt}}$ тогда становится мерой того, насколько система может вырасти, прежде чем она должна будет разделиться или реорганизоваться. Тот факт, что $k$ имеет порядок 1,5–2 в разных системах, предполагает, что лежащее в основе отношение $K_{\mathrm{eff},\text{crit}}/K_{\mathrm{eff},\text{opt}}$ универсально, что согласуется с фрактальным законом масштабирования. Таким образом, представленные здесь данные косвенно поддерживают существование технологического порога и его связь с универсальным показателем $\beta$.
	
	\section{Заключение}
	\label{sec:conclusion}
	
	Собранные в этом приложении данные, хотя и предварительные, позволяют предположить, что критические размеры групп, при которых коллективные системы претерпевают фазовые переходы, не являются произвольными, а следуют универсальному закону масштабирования, связанному с показателем YUCT $\beta = 0,67$ и аксиомой технологического словарного порога. Отношение $k = N_{\text{crit}}/N_{\text{opt}}$ группируется вокруг 1,5–2,5, и вытекающий из него порог ошибки $\varepsilon_{\text{crit}}/\varepsilon_{\text{opt}} \approx 1,3$–1,6 является удивительно согласованным для разных систем. Мы наметили конкретные проверяемые предсказания, которые могут быть проверены с помощью будущих экспериментов и наблюдений. Подтверждение этих предсказаний обеспечило бы сильную независимую поддержку YUCT, а также фундаментальной природы метрики эффективности координации $\Keff$ и её законов масштабирования.
	
	\begin{thebibliography}{99}
		\bibitem{Seeley2010} Seeley, T. D. (2010). \emph{Honeybee Democracy}. Princeton University Press.
		\bibitem{Winston1987} Winston, M. L. (1987). \emph{The Biology of the Honey Bee}. Harvard University Press.
		\bibitem{FM3-90.5} FM 3-90.5 (2008). \emph{Combined Arms Battalion}. US Army Field Manual.
		\bibitem{Tunstrom2013} Tunstrøm, K. et al. (2013). Collective states, multistability and transitional behavior in schooling fish. \emph{PLoS Computational Biology}, 9(2), e1002915.
		\bibitem{Connor2000} Connor, R. C. et al. (2000). The bottlenose dolphin: social relationships in a fission‑fusion society. In: \emph{Cetacean Societies}, University of Chicago Press.
		\bibitem{FAO2013} FAO (2013). \emph{In vivo conservation of animal genetic resources}. FAO Animal Production and Health Guidelines No. 14.
		\bibitem{Alberts1995} Alberts, S. C. \& Altmann, J. (1995). Balancing costs and opportunities: dispersal in male baboons. \emph{American Naturalist}, 145(2), 279–306.
		\bibitem{Yakushev2026} Yakushev, A. V. (2026). \emph{Yakushev Unified Coordination Theory (YUCT)}. Zenodo. \url{https://doi.org/10.5281/zenodo.18444599}
		\bibitem{AppendixL} Приложение L: Фрактальное масштабирование ошибки координации – универсальный закон от ДНК до космологии.
		\bibitem{AppendixO} Приложение O: Аксиома технологического словарного порога.
		% Добавлена недостающая ссылка для \cite{bees}
		\bibitem{bees} Seeley, T. D. (2010) – см. выше, или заменить на конкретную ссылку.
	\end{thebibliography}
	
\end{document}